\chapter{Introduction}
\label{chap:introduction}

\section{Background}

Large language models can generate fluent explanations and procedural
instructions, but their parameters are not an auditable or easily updated
repository of organizational knowledge. Retrieval-Augmented Generation (RAG)
addresses this limitation by combining parametric generation with a
non-parametric document collection \cite{lewis2020rag}. Instead of relying only
on facts remembered during pre-training, a RAG system searches an external
knowledge base, places relevant passages in the model prompt, and asks the model
to produce an answer grounded in those passages.

This pattern is attractive for enterprise question answering because policies,
work instructions, support articles, and operational procedures change more
frequently than foundation models can be retrained. Retrieval also provides
provenance: a user can inspect which source was used and an engineer can examine
whether failure originated in retrieval or generation. These properties are
particularly important for business workflows, where an answer must often
preserve prerequisites, ordered steps, conditional branches, exceptions, and
responsibility boundaries.

However, the conventional RAG pipeline is usually static. Every question is
processed with the same search method, the same number of retrieved chunks, and
the same generation route. A direct definition may therefore incur unnecessary
retrieval and model cost, while a multi-stage workflow question may receive only
one retrieval pass and miss critical evidence. Adaptive-RAG addresses this
problem by predicting query complexity with a smaller language model and
selecting among no retrieval, single-step retrieval, and iterative retrieval
\cite{jeong2024adaptiverag}. The central principle is to use the least expensive
strategy that is sufficiently capable for the question.

The suitability of an adaptive approach for enterprise QA cannot be established
using implementation functionality alone. It requires a benchmark whose
questions and answers are grounded in a fixed enterprise knowledge base. WixQA
was created for this purpose and contains expert-written, simulated, and
synthetic QA datasets aligned to a Wix Help Center corpus
\cite{cohen2025wixqa}. Evaluation also needs to be actionable. Aggregate
faithfulness or relevance scores indicate whether a system performs poorly, but
do not necessarily identify whether the remedy is better chunking, a larger
\textit{top-k}, reranking, prompt changes, or another route. RAGXplain addresses
that gap by converting metric evidence into explanatory findings and
recommendations \cite{cohen2025ragxplain}.

\section{Problem Statement}

Enterprise workflow QA creates four related engineering and research problems.
First, incoming questions have different computational needs. A fixed route
either wastes resources on easy questions or under-serves complex questions.
Second, retrieval quality depends on a chain of decisions--document loading,
deduplication, chunking, embedding, lexical indexing, vector compatibility,
filtering, and reranking. An error at any point can remove the evidence needed
by the generator. Third, model selection introduces operational concerns such
as credentials, data egress, cost, provider availability, streaming, and
fallback behavior. Fourth, a multi-stage answer is difficult to trust unless
its sources, route, model calls, timings, and failures can be reconstructed.

The problem addressed by this project is therefore:

\begin{quote}
How can a modular Adaptive RAG system classify the complexity of business
workflow questions, select a proportionate execution strategy, produce
source-grounded answers, and provide sufficient evaluation and observability
evidence to improve the system?
\end{quote}

The work is not limited to a classifier experiment or a RAG chatbot mock-up. The
classifier, knowledge pipeline, retrieval strategies, model gateway,
evaluation workflow, and user interface must operate as one reproducible
system. At the same time, research results must be reported conservatively due to resource constrainst:
the completed 20-case compatible experiment establishes current evidence, but
it cannot substitute for a larger controlled benchmark matrix.

\section{Research Objectives}

The research objectives from the proposal were retained and refined as the
implementation matured. \Cref{tab:research-objectives} links each objective to
its system evidence.

\begin{table}[H]
  \centering
  \caption{Research objectives and implementation evidence.}
  \label{tab:research-objectives}
  \begin{tabularx}{\textwidth}{>{\bfseries}p{0.11\textwidth}X X}
    \toprule
    Objective & Research intent & Evidence produced by the project\\
    \midrule
    RO1 &
    Develop a lightweight domain-specialized query-complexity classifier for
    adaptive routing. &
    Reproducible QAC preparation, a balanced four-class dataset, trained
    DistilBERT and T5-small artifacts, calibrated routing interfaces, and
    grouped-validation comparison metrics.\\
    RO2 &
    Design and implement an Adaptive RAG framework that selects an execution
    path appropriate to query complexity. &
    L1 direct generation, L2 single-step RAG, L3 decomposition and aggregated
    retrieval, and bounded L4 planning with traceable fallback.\\
    RO3 &
    Evaluate answer quality, retrieval, routing, efficiency, and failure
    causes. &
    WixQA experiments, lexical and retrieval metrics, Model Farm judge calls,
    RAGXplain diagnostics, trace artifacts, cost/latency telemetry, and user
    feedback analytics.\\
    \bottomrule
  \end{tabularx}
\end{table}

\section{Research Questions}

\begin{description}[style=nextline,leftmargin=1.5cm]
  \item[RQ1]
  How accurately and efficiently can a lightweight, domain-specialized model
  classify business workflow query complexity and route questions among direct,
  single-step, multi-step, and advanced execution?

  \item[RQ2]
  Does the Adaptive RAG framework improve answer faithfulness and groundedness
  while controlling latency, token consumption, and cost compared with fixed
  configurations?

  \item[RQ3]
  Can explainable evaluation through RAGXplain turn metric outputs into
  actionable recommendations for retrieval, generation, and configuration
  improvement?
\end{description}

The original Adaptive-RAG study defines three strategies
\cite{jeong2024adaptiverag}. The \texttt{advanced} class and L4 route in RQ1
are a project extension motivated by questions that require bounded planning or
tool use. They are not represented as an original contribution of that paper.

\section{Scope}

\subsection{Included Scope}

The data acquisition and preprocessing scope covers the three official WixQA
question-answering sources and the associated business-workflow knowledge
corpus. The system normalizes these data into reproducible question--answer--
context records and supports ingestion from local TXT, JSON, JSONL, Markdown,
PDF, and DOCX files, the prepared WixQA corpus, and security-checked public web
pages. Knowledge processing includes metadata extraction, content-hash
deduplication, configurable and structure-aware chunking, embedding, and
atomic activation of versioned indexes in PostgreSQL and pgVector.

The model-specialization scope includes constructing a balanced four-class
query-complexity dataset and training lightweight DistilBERT and T5-small
classifiers. These models categorize questions as \textit{simple},
\textit{moderate}, \textit{complex}, or \textit{advanced} and provide the
evidence used by Adaptive routing. Foundation and hosted models are integrated
through a provider-neutral Model Gateway rather than trained from scratch. The
gateway supports local models and selected LiteLLM-compatible providers,
encrypted credentials, deployment health checks, budget controls, and usage
telemetry.

The Adaptive RAG scope covers four bounded execution paths: L1 direct
generation without retrieval, L2 single-step retrieval, L3 decomposed
multi-query retrieval, and L4 iterative planning with authorized tools. The
retrieval layer supports BM25, dense, and hybrid search, document filtering,
and optional reranking. Query-time processing also includes conversation-aware
reformulation, streaming responses, source citations, cancellation, answer
regeneration, retry, and full-fidelity execution tracing.

The prototype and evaluation scope consists of a React-based AI assistant and
a FastAPI service backed by PostgreSQL. The application provides persistent
conversation history, reusable RAG configurations, Knowledge Base and AI Model
administration, feedback, and usage analytics. It evaluates saved
configurations against compatible WixQA cases and integrates RAGXplain to
translate semantic metrics into explainable findings and actionable
configuration recommendations. This prototype supports iterative refinement
and controlled comparison with fixed RAG routes, while remaining a capstone
research platform rather than a production service.

\subsection{Excluded or Deferred Scope}

The project does not target general-purpose open-domain reasoning, real-time
audio or video processing, or unrestricted multimodal interaction. It does not
pre-train a foundation model from scratch; instead, it specializes existing
pre-trained classifiers and consumes existing local or hosted generation,
embedding, reranking, planning, and judging models.

Large-scale decentralized multi-agent architectures are also excluded. L4 is
implemented as a controlled agentic extension with explicit limits on planner
iterations, tool calls, execution time, cost, data egress, and cancellation.
Google Drive, OneDrive, and general database connectors remain unavailable
extension points rather than production integrations. Production
multi-tenancy, enterprise identity federation, high-availability deployment,
managed disaster recovery, and formal penetration testing are deferred beyond
the capstone.

\section{Contributions}

The principal contributions are:

\begin{enumerate}
  \item a full-stack Adaptive RAG reference architecture that keeps classifier,
  retrieval, generation, evaluation, and observability interfaces separable;
  \item a four-level routing model that preserves the least-complex-successful
  principle and adds a bounded agentic extension;
  \item a versioned knowledge-processing pipeline that prevents mixed embedding
  indexes and supports a token-aware WixQA/MiniLM chunking profile;
  \item a unified AI Model Gateway for in-process, hosted, direct-provider, and
  self-hosted model deployments;
  \item an evaluation workflow that treats each saved RAG configuration as a
  reproducible experimental condition;
  \item full-fidelity hierarchical traces that connect route decisions,
  retrieval diagnostics, model usage, sources, costs, and answer versions; and
  \item a practical studio interface for configuring, operating, evaluating,
  and diagnosing the system.
\end{enumerate}

\section{Report Organization}

\Cref{chap:background} reviews the research and technology foundations.
\Cref{chap:methodology} defines the data, routing, knowledge-processing, and
evaluation methods. \Cref{chap:implementation-results} describes the
implementation and reports verified or preliminary evidence.
\Cref{chap:discussion} discusses the research questions, limitations,
conclusions, and future work. The appendices provide setup instructions,
schemas, API and persistence summaries, detailed experiment references, and
evidence placeholders.
